\exercisesheader{}

% 1

\eoce{\qt{Tentukan nilai kritis $t$\label{identify_critical_t}} Suatu sampel acak
independen dipilih dari populasi yang kira-kira berdistribusi normal dengan
simpangan baku tidak diketahui. Tentukan derajat kebebasan dan nilai kritis $t$
(t$^\star$) untuk ukuran sampel dan tingkat kepercayaan yang diberikan.
%\begin{multicols}{4}
\begin{parts}
\item $n = 6$, TK = 90\%
\item $n = 21$, TK = 98\%
\item $n = 29$, TK = 95\%
\item $n = 12$, TK = 99\%
\end{parts}
%\end{multicols}
}{}

% 2

\eoce{\qt{Distribusi $t$\label{t_distribution}}
Gambar di sebelah kanan memperlihatkan tiga kurva
unimodal dan simetris:
distribusi normal baku (z),
distribusi $t$ dengan 5 derajat kebebasan,
dan distribusi $t$ dengan 1 derajat kebebasan.
Tentukan kurva mana yang mewakili masing-masing distribusi, lalu jelaskan alasan Anda.
\begin{center}

\FigureFullPath[Tiga distribusi ditampilkan; semuanya simetris, berbentuk lonceng, dan berpusat di nol. Kurva pertama digambar dengan garis utuh dan memiliki puncak tengah paling tinggi dan tajam; ekornya paling tipis dan tampak mendekati nol sekitar -3 dan 3. Kurva kedua digambar dengan garis putus-putus dan memiliki tinggi puncak serta ketebalan ekor di antara kedua kurva lain; ekornya mendekati nol sekitar -4 dan 4. Kurva ketiga digambar dengan garis titik-titik dan memiliki puncak tengah paling rendah dan lebar; ekornya paling tebal dan masih tampak melampaui -4 dan 4.]{0.4}{ch_inference_for_means/figures/eoce/t_distribution/t_distribution}
\end{center}
}{}

% 3

\eoce{\qt{Tentukan nilai-p, Bagian I\label{find_T_pval_1_2_sided}}
Suatu sampel acak independen
dipilih dari populasi yang kira-kira berdistribusi normal
dengan simpangan baku
yang tidak diketahui.
Tentukan nilai-p untuk ukuran sampel dan statistik uji yang diberikan.
Tentukan juga apakah hipotesis nol akan ditolak pada
$\alpha = 0.05$.
\begin{parts}
\item $n = 11$, $T = 1.91$
\item $n = 17$, $T = -3.45$
\item $n = 7$, $T = 0.83$
\item $n = 28$, $T = 2.13$
\end{parts}
}{}

% 4

\eoce{\qt{Tentukan nilai-p, Bagian II\label{find_T_pval_2_2_sided}}
Suatu sampel acak independen
dipilih dari populasi yang kira-kira berdistribusi normal
dengan simpangan baku
yang tidak diketahui.
Tentukan nilai-p untuk ukuran sampel dan statistik uji yang diberikan.
Tentukan juga apakah hipotesis nol akan ditolak pada
$\alpha = 0.01$.
\begin{parts}
\item $n = 26$, $T = 2.485$
\item $n = 18$, $T = 0.5$
\end{parts}
}{}

% 5

\eoce{\qt{Menelusuri perhitungan secara terbalik, Bagian I\label{work_backwards_1}} Interval kepercayaan 95\%
untuk rata-rata populasi, $\mu$, diberikan sebagai (18.985, 21.015). Interval
kepercayaan ini didasarkan pada sampel acak sederhana berisi 36 pengamatan.
Hitung rata-rata dan simpangan baku sampel. Asumsikan bahwa semua kondisi
yang diperlukan untuk inferensi terpenuhi. Gunakan distribusi $t$ dalam setiap
perhitungan.
}{}

% 6

\eoce{\qt{Menelusuri perhitungan secara terbalik, Bagian II\label{work_backwards_2}} Interval kepercayaan 90\%
untuk rata-rata populasi adalah (65, 77). Distribusi populasi kira-kira
normal dan simpangan baku populasi tidak diketahui. Interval
kepercayaan ini didasarkan pada sampel acak sederhana berisi 25 pengamatan.
Hitung rata-rata sampel, margin galat, dan simpangan baku
sampel.
}{}

\D{\newpage}

% 7

\eoce{\qt{Kebiasaan tidur warga New York\label{ny_sleep_habits_2_sided}}
New York dikenal sebagai
``kota yang tidak pernah tidur''.
Suatu sampel acak yang terdiri atas 25 warga New York ditanya berapa lama
mereka tidur setiap malam.
Ringkasan statistik data tersebut ditampilkan
di bawah ini.
Estimasi titik menunjukkan bahwa, secara rata-rata, warga New York tidur kurang dari
8~jam per malam.
Apakah hasilnya signifikan secara statistik?
\begin{center}
\begin{tabular}{rrrrrr}
 \hline
 n   & $\bar{x}$ & s     & min   & maks \\
 \hline
 25  & 7.73      & 0.77  & 6.17  & 9.78 \\
  \hline
\end{tabular}
\end{center}

\begin{parts}
\item Tuliskan hipotesis dalam simbol dan kata-kata.
\item Periksa kondisi, lalu hitung statistik uji, $T$, dan
derajat kebebasan yang terkait.
\item Tentukan dan tafsirkan nilai-p dalam konteks ini. Menggambar sketsa mungkin
membantu.
\item Apa kesimpulan uji hipotesis tersebut?
\item Jika Anda menyusun interval kepercayaan 90\% yang bersesuaian dengan
uji hipotesis ini, apakah Anda memperkirakan 8 jam akan berada dalam interval tersebut?
\end{parts}
}{}

% 8

\eoce{\qt{Tinggi badan orang dewasa\label{adult_heights}}
Peneliti yang mengkaji antropometri
mengumpulkan ukuran lingkar tubuh dan diameter rangka, serta
usia, berat badan, tinggi badan, dan jenis kelamin, dari 507 individu yang aktif secara fisik. Histogram
di bawah ini menunjukkan distribusi sampel tinggi badan dalam sentimeter.
\footfullcite{Heinz:2003} \\
\begin{minipage}[c]{0.75\textwidth}
\begin{center}
\FigureFullPath[Histogram ditampilkan untuk "Tinggi badan" dengan nilai dari 140 sampai 200 dan lebar bin 5. Distribusinya kira-kira simetris dengan pusat sekitar 170. Tinggi bin, dimulai dari bin 145 sampai 150, kira-kira 3, 17, 55, 70, 100, 85, 95, 50, 30, 15, dan 3.]{}{ch_inference_for_means/figures/eoce/adult_heights/adult_heights_hist}
\end{center}
\end{minipage}
\begin{minipage}[c]{0.23\textwidth}
\begin{center}
\begin{tabular}{l|r}
Min     & 147.2 \\
Q1      & 163.8 \\
Median  & 170.3 \\
Rata-rata & 171.1 \\
SD      &  9.4 \\
Q3      & 177.8 \\
Maks    & 198.1 \\
\end{tabular}
\end{center}
\end{minipage}
\begin{parts}
\item Berapakah estimasi titik untuk rata-rata tinggi badan individu yang aktif?
Bagaimana dengan mediannya?
\item Berapakah estimasi titik untuk simpangan baku tinggi badan
individu yang aktif? Bagaimana dengan IQR?
\item Apakah orang dengan tinggi 1m 80cm (180 cm) dianggap luar biasa tinggi? Dan apakah
orang dengan tinggi 1m 55cm (155cm) dianggap luar biasa pendek? Jelaskan
alasan Anda.
\item Para peneliti mengambil sampel acak lain dari individu yang aktif secara fisik.
Apakah Anda memperkirakan rata-rata dan simpangan baku sampel baru ini
akan sama dengan nilai yang diberikan di atas? Jelaskan alasan Anda.
\item Rata-rata sampel yang diperoleh merupakan estimasi titik bagi rata-rata tinggi badan semua
individu yang aktif, jika sampel individu tersebut setara dengan sampel
acak sederhana.
Ukuran apa yang kita gunakan untuk menguantifikasi variabilitas estimasi semacam itu?
Hitung
besaran tersebut menggunakan data dari sampel awal dengan syarat bahwa
data itu merupakan sampel acak sederhana.
\end{parts}
}{}

% 9

\eoce{\qt{Tentukan rata-rata\label{find_mean_2_sided}}
Anda diberi hipotesis berikut:
\begin{align*}
H_0&: \mu = 60 \\
H_A&: \mu \neq 60
\end{align*}
Kita mengetahui bahwa simpangan baku sampel adalah 8
dan ukuran sampelnya 20.
Untuk nilai rata-rata sampel berapakah nilai-p sama dengan 0.05?
Asumsikan bahwa semua kondisi yang diperlukan untuk inferensi terpenuhi.
}{}

\D{\newpage}

% 10

\eoce{\qt{$t^\star$ versus $z^\star$\label{critical_t_vs_z}} Untuk tingkat kepercayaan
tertentu, $t^{\star}_{df}$ lebih besar daripada $z^{\star}$. Jelaskan bagaimana nilai $t^{\star}_{df}$
yang sedikit lebih besar daripada $z^{\star}$ memengaruhi lebar interval kepercayaan.
}{}

% 11

\eoce{\qt{Bermain piano\label{play_piano_2_sided}}
Georgianna menyatakan bahwa di sebuah kota kecil
yang terkenal karena sekolah musiknya, rata-rata anak mengikuti les piano kurang dari 5 tahun.
Kita memiliki sampel acak berisi 20 anak dari kota tersebut, dengan
rata-rata lama les piano 4.6 tahun dan simpangan baku 2.2 tahun.
\begin{parts}
\item
    Evaluasi pernyataan Georgianna (atau kemungkinan bahwa kebalikannya benar)
    menggunakan uji hipotesis.
\item
    Susun interval kepercayaan 95\% untuk banyaknya tahun
    siswa di kota ini mengikuti les piano, lalu tafsirkan interval itu
    dalam konteks data.
\item
    Apakah hasil uji hipotesis dan interval
    kepercayaan Anda konsisten?
    Jelaskan alasan Anda.
\end{parts}
}{}

% 12

\eoce{\qt{Gas buang kendaraan dan
    paparan timbal\label{auto_exhaust_lead_exposure_2_sided}}
Peneliti yang tertarik pada paparan timbal akibat gas buang kendaraan
mengambil sampel darah dari 52 petugas polisi yang terus-menerus
menghirup asap knalpot kendaraan ketika bertugas mengatur lalu lintas
di lingkungan yang terutama bersifat perkotaan.
Sampel darah para petugas tersebut memiliki rata-rata konsentrasi timbal
124.32 $\mu$g/l dan SD 37.74 $\mu$g/l;
penelitian sebelumnya terhadap individu dari daerah pinggiran kota terdekat,
yang tidak memiliki riwayat paparan, menemukan rata-rata kadar
timbal darah
sebesar 35 $\mu$g/l.\footfullcite{Mortada:2000}
\begin{parts}
\item
    Tuliskan hipotesis yang sesuai untuk
    menguji apakah para petugas polisi tampaknya telah terpapar
    konsentrasi timbal yang berbeda.
\item\label{auto_exhaust_lead_exposure_2_sided_cond}
    Nyatakan secara eksplisit dan periksa semua kondisi yang diperlukan untuk
    melakukan inferensi pada data ini.
\item
    Terlepas dari jawaban Anda pada
    bagian~(\ref{auto_exhaust_lead_exposure_2_sided_cond}),
    ujilah hipotesis bahwa petugas polisi pusat kota memiliki
    paparan timbal yang lebih tinggi daripada kelompok dalam penelitian sebelumnya.
    Tafsirkan hasil Anda dalam konteksnya.
\end{parts}
}{}

% 13

\eoce{\qt{Penghematan asuransi mobil\label{car_insurance_savings}}
Seorang peneliti pasar ingin mengevaluasi penghematan asuransi mobil
di perusahaan pesaing.
Berdasarkan penelitian sebelumnya, ia mengasumsikan bahwa simpangan
baku penghematan adalah \$100.
Ia ingin mengumpulkan data sehingga memperoleh margin
galat tidak lebih dari \$10 pada tingkat kepercayaan 95\%.
Seberapa besar sampel yang harus ia kumpulkan?
}{}

% 14

\eoce{\qt{Skor SAT\label{sat_scores_CI}}
Simpangan baku skor SAT mahasiswa di
sebuah perguruan tinggi Ivy League tertentu adalah 250 poin.
Dua mahasiswa statistika, Raina dan Luke, ingin mengestimasi
rata-rata skor SAT mahasiswa di perguruan tinggi ini sebagai bagian
dari proyek kelas.
Mereka ingin margin galatnya tidak lebih dari 25 poin.
\begin{parts}
\item
    Raina ingin menggunakan interval kepercayaan 90\%.
    Seberapa besar sampel yang harus ia kumpulkan?
\item
    Luke ingin menggunakan interval kepercayaan 99\%.
    Tanpa menghitung ukuran sampel yang sebenarnya, tentukan
    apakah sampelnya harus lebih besar atau lebih kecil
    daripada sampel Raina, lalu jelaskan alasan Anda.
\item
    Hitung ukuran sampel minimum yang diperlukan Luke.
\end{parts}
}{}
